% ============================================================
%  SCHOLA DOMESTICA — Magister Mathematica
%  Level 1 | Week 5 | Day 1
%  Topic: Numbers 6–10 — More, Fewer, or the Same?
%  Enrichment: Roman Numerals VI–X
% ============================================================
\documentclass[12pt, letterpaper]{article}

% ---------- packages ----------
\usepackage[top=1in, bottom=1in, left=1in, right=1in]{geometry}
\usepackage{fontenc}
\usepackage{inputenc}
\usepackage{lmodern}
\usepackage{microtype}
\usepackage{amsmath}
\usepackage{amssymb}
\usepackage{graphicx}
\usepackage{xcolor}
\usepackage{tikz}
\usepackage{array}
\usepackage{booktabs}
\usepackage{enumitem}
\usepackage{multicol}
\usepackage{mdframed}
\usepackage{fancyhdr}
\usepackage{calc}

% ---------- colours ----------
\definecolor{scholared}{RGB}{160,30,30}
\definecolor{scholagold}{RGB}{180,140,40}
\definecolor{scholacream}{RGB}{253,248,235}
\definecolor{scholagreen}{RGB}{50,110,60}
\definecolor{scholablue}{RGB}{40,80,140}

% ---------- page style ----------
\pagestyle{fancy}
\fancyhf{}
\renewcommand{\headrulewidth}{0.6pt}
\fancyhead[L]{\small\textcolor{scholared}{\textbf{Schola Domestica — Magister Mathematica}}}
\fancyhead[R]{\small\textcolor{scholared}{Level 1 · Week 5 · Day 1}}
\fancyfoot[C]{\small\textcolor{gray}{— \thepage\ —}}

% ---------- custom commands ----------
\newcommand{\dotbox}[1]{%
  \tikz[baseline=-0.5ex]{\draw[scholablue, thick] (0,0) rectangle (#1 cm, 0.9cm);}%
}
\newcommand{\answerline}{\underline{\hspace{2.5cm}}}
\newcommand{\biganswerline}{\underline{\hspace{4cm}}}

% draw a row of circles (filled) to represent objects
\newcommand{\dots}[1]{%
  \begin{tikzpicture}[baseline=-0.5ex]
    \foreach \i in {1,...,#1} {
      \filldraw[scholablue] (\i*0.45-0.4, 0) circle (0.15cm);
    }
  \end{tikzpicture}%
}

% ---------- section box ----------
\newmdenv[
  backgroundcolor=scholacream,
  linecolor=scholared,
  linewidth=1.5pt,
  roundcorner=6pt,
  innertopmargin=8pt,
  innerbottommargin=8pt,
  innerleftmargin=10pt,
  innerrightmargin=10pt
]{lessonbox}

\begin{document}

% ===== TITLE =====
\begin{center}
  {\LARGE\textbf{\textcolor{scholared}{Numbers 6–10:}}}\\[4pt]
  {\Large\textcolor{scholagold}{\textit{More, Fewer, or the Same?}}}\\[6pt]
  {\small Name: \biganswerline \hspace{1cm} Date: \answerline}
\end{center}

\vspace{4pt}
\noindent\textcolor{scholared}{\rule{\linewidth}{1pt}}
\vspace{4pt}

% ===== PICTURE INSTRUCTION =====
% IMAGE PROMPT (Beatrix Potter watercolour style):
% Two groups of animals facing each other: on the left, 7 small brown field mice sitting upright;
% on the right, 4 golden acorns arranged in a row. Soft watercolour wash, warm autumn palette,
% white paper background, no text. Proportions suitable for a half-page illustration.

\begin{lessonbox}
\textbf{Remember:}\\[4pt]
\begin{itemize}[leftmargin=1.5em, itemsep=2pt]
  \item \textbf{More} means a bigger group. 
  \item \textbf{Fewer} means a smaller group.
  \item \textbf{Same} means both groups are equal.
\end{itemize}
\medskip
\noindent To compare, \emph{match one object from each group}. The group with leftover objects has \textbf{more}.
\end{lessonbox}

\bigskip

% ===== SECTION 1: COUNT & COMPARE =====
\noindent\textcolor{scholared}{\large\textbf{Part I — Count and Compare}} \hfill \textit{(Problems 1–4)}

\medskip
\noindent\textit{Count each group. Write the number. Circle the word.}

\bigskip

% Problem 1
\noindent
\begin{tabular}{@{}p{5.5cm}@{\hspace{0.5cm}}p{5.5cm}@{}}
  \textbf{1.} Left group: \quad
  \begin{tikzpicture}[baseline=-0.5ex]
    \foreach \i in {1,...,6} { \filldraw[scholablue] (\i*0.38, 0) circle (0.13cm); }
  \end{tikzpicture}
  &
  Right group: \quad
  \begin{tikzpicture}[baseline=-0.5ex]
    \foreach \i in {1,...,8} { \filldraw[scholagreen] (\i*0.38, 0) circle (0.13cm); }
  \end{tikzpicture}
\end{tabular}

\medskip
\hspace{1.5em}Left: \answerline \quad Right: \answerline

\medskip
\hspace{1.5em}The left group has \quad \textbf{MORE} \quad / \quad \textbf{FEWER} \quad / \quad \textbf{THE SAME} \quad as the right group.

\bigskip\bigskip

% Problem 2
\noindent
\begin{tabular}{@{}p{5.5cm}@{\hspace{0.5cm}}p{5.5cm}@{}}
  \textbf{2.} Left group: \quad
  \begin{tikzpicture}[baseline=-0.5ex]
    \foreach \i in {1,...,7} { \filldraw[scholared] (\i*0.38, 0) circle (0.13cm); }
  \end{tikzpicture}
  &
  Right group: \quad
  \begin{tikzpicture}[baseline=-0.5ex]
    \foreach \i in {1,...,7} { \filldraw[scholagold] (\i*0.38, 0) circle (0.13cm); }
  \end{tikzpicture}
\end{tabular}

\medskip
\hspace{1.5em}Left: \answerline \quad Right: \answerline

\medskip
\hspace{1.5em}The left group has \quad \textbf{MORE} \quad / \quad \textbf{FEWER} \quad / \quad \textbf{THE SAME} \quad as the right group.

\bigskip\bigskip

% Problem 3
\noindent
\begin{tabular}{@{}p{5.5cm}@{\hspace{0.5cm}}p{5.5cm}@{}}
  \textbf{3.} Left group: \quad
  \begin{tikzpicture}[baseline=-0.5ex]
    \foreach \i in {1,...,9} { \filldraw[scholablue] (\i*0.38, 0) circle (0.13cm); }
  \end{tikzpicture}
  &
  Right group: \quad
  \begin{tikzpicture}[baseline=-0.5ex]
    \foreach \i in {1,...,6} { \filldraw[scholagreen] (\i*0.38, 0) circle (0.13cm); }
  \end{tikzpicture}
\end{tabular}

\medskip
\hspace{1.5em}Left: \answerline \quad Right: \answerline

\medskip
\hspace{1.5em}The left group has \quad \textbf{MORE} \quad / \quad \textbf{FEWER} \quad / \quad \textbf{THE SAME} \quad as the right group.

\bigskip\bigskip

% Problem 4
\noindent
\begin{tabular}{@{}p{5.5cm}@{\hspace{0.5cm}}p{5.5cm}@{}}
  \textbf{4.} Left group: \quad
  \begin{tikzpicture}[baseline=-0.5ex]
    \foreach \i in {1,...,10} { \filldraw[scholared] (\i*0.38, 0) circle (0.13cm); }
  \end{tikzpicture}
  &
  Right group: \quad
  \begin{tikzpicture}[baseline=-0.5ex]
    \foreach \i in {1,...,8} { \filldraw[scholagold] (\i*0.38, 0) circle (0.13cm); }
  \end{tikzpicture}
\end{tabular}

\medskip
\hspace{1.5em}Left: \answerline \quad Right: \answerline

\medskip
\hspace{1.5em}The left group has \quad \textbf{MORE} \quad / \quad \textbf{FEWER} \quad / \quad \textbf{THE SAME} \quad as the right group.

\bigskip
\noindent\textcolor{scholared}{\rule{\linewidth}{0.4pt}}

% ===== SECTION 2: ORDERING =====
\noindent\textcolor{scholared}{\large\textbf{Part II — Put in Order}} \hfill \textit{(Problems 5–8)}

\medskip
\noindent\textit{Write these numbers from \textbf{least to greatest} (smallest to largest).}

\bigskip
\noindent\textbf{5.} \quad 9 \quad 6 \quad 8 \quad 7 \hspace{2cm}
  \answerline \; $<$ \; \answerline \; $<$ \; \answerline \; $<$ \; \answerline

\bigskip
\noindent\textbf{6.} \quad 10 \quad 6 \quad 9 \quad 7 \hspace{1.5cm}
  \answerline \; $<$ \; \answerline \; $<$ \; \answerline \; $<$ \; \answerline

\bigskip
\noindent\textbf{7.} \quad 8 \quad 10 \quad 6 \quad 9 \hspace{1.5cm}
  \answerline \; $<$ \; \answerline \; $<$ \; \answerline \; $<$ \; \answerline

\bigskip
\noindent\textit{Write these numbers from \textbf{greatest to least} (largest to smallest).}

\bigskip
\noindent\textbf{8.} \quad 6 \quad 10 \quad 8 \quad 7 \hspace{1.5cm}
  \answerline \; $>$ \; \answerline \; $>$ \; \answerline \; $>$ \; \answerline

\bigskip
\noindent\textcolor{scholared}{\rule{\linewidth}{0.4pt}}

% ===== ENRICHMENT =====
\begin{lessonbox}
\textcolor{scholared}{\textbf{Enrichment — Roman Numerals VI through X}}\\[6pt]
The Romans wrote 6–10 like this:\\[6pt]
\begin{center}
\renewcommand{\arraystretch}{1.5}
\begin{tabular}{|c|c|c|c|c|}
\hline
\textbf{VI} & \textbf{VII} & \textbf{VIII} & \textbf{IX} & \textbf{X} \\
\hline
6 & 7 & 8 & 9 & 10 \\
\hline
\end{tabular}
\end{center}
\medskip
\noindent \textit{Trace and copy each Roman numeral once:}

\medskip
VI \quad \dotbox{1.5} \hspace{0.5cm}
VII \quad \dotbox{1.5} \hspace{0.5cm}
VIII \quad \dotbox{1.5} \hspace{0.5cm}
IX \quad \dotbox{1.5} \hspace{0.5cm}
X \quad \dotbox{1.5}

\medskip
\noindent\textit{What Roman numeral means 8?} \quad \answerline \qquad
\textit{What does IX mean?} \quad \answerline
\end{lessonbox}

\end{document}
